\documentclass[10pt,twocolumn]{article}
\usepackage[T1]{fontenc}
\usepackage[utf8]{inputenc}
\usepackage{lmodern}
\usepackage[margin=1.8cm,top=2.4cm,bottom=2cm,columnsep=0.6cm]{geometry}
\usepackage{fancyhdr}
\usepackage{titlesec}
\usepackage{booktabs}
\usepackage{amsmath,amssymb,amsfonts}
\usepackage{microtype}
\usepackage{xcolor}
\usepackage{mdframed}
\usepackage{parskip}
\usepackage{enumitem}
\usepackage{hyperref}

\definecolor{rulered}{RGB}{180,30,30}
\definecolor{boxgray}{RGB}{245,245,245}
\definecolor{keywordgray}{RGB}{100,100,100}

\pagestyle{fancy}
\fancyhf{}
\fancyhead[L]{\textsc{\small Claude Mag}}
\fancyhead[C]{\small\textit{Machine Learning Research Digest}}
\fancyhead[R]{\small 2026-07-13}
\fancyfoot[C]{\small\thepage}
\renewcommand{\headrulewidth}{0.8pt}

\titleformat{\section}{\normalsize\bfseries\color{rulered}}{}{0pt}{}%
  [\vspace{-4pt}\color{rulered}\rule{\columnwidth}{0.4pt}]
\titlespacing{\section}{0pt}{8pt}{4pt}

\hypersetup{colorlinks=true,urlcolor=rulered,linkcolor=black}

\begin{document}
\setlength{\parindent}{0pt}
\setlength{\parskip}{4pt}

{\small\color{keywordgray}\textbf{Keywords:}
  model compression \textbullet{}
  binary quantization \textbullet{}
  large language models \textbullet{}
  spherical coding}\par\vspace{4pt}

{\LARGE\bfseries\raggedright
  BiSCo-LLM}\par\vspace{4pt}

{\large\itshape\raggedright
  Lookup-Free Binary Spherical Coding Brings
  Sub-1-Bit LLM Compression Within Reach}\par\vspace{6pt}

{\small\textbf{By} Yuantian Shao, Peisong Wang, Zhilei Liu,
  Chuangyi Li, Yuanteng Chen, Pengcheng Xie, Yiwu Yao, Zhihui Wei \& Jian Cheng
  (Chinese Academy of Sciences; University of CAS; Huawei Technologies)}\par\vspace{2pt}

{\small\color{keywordgray}arXiv:2607.08643 \textbullet{}
  arXiv preprint, July 2026}
\par\vspace{2pt}\noindent{\color{rulered}\rule{\columnwidth}{0.6pt}}\vspace{6pt}

Large language models are memory-bound: frontier models require tens of
gigabytes of weight storage, creating an insurmountable barrier for
on-device and edge deployment. Existing extreme-quantization methods
face a fundamental trade-off---codebook-based approaches achieve good
quality but require storing a codebook and paying a lookup cost at
inference; binary approaches eliminate lookup overhead but suffer severe
quality degradation. BiSCo-LLM resolves this trade-off by introducing
Binary Spherical Coding (BSCo), a lookup-table-free scheme that matches
2-bit codebook quality at 1-bit storage by representing weights as
binary combinations of spherical frame vectors.

\begin{mdframed}[backgroundcolor=boxgray,linecolor=rulered,linewidth=1pt,
    innerleftmargin=6pt,innerrightmargin=6pt,
    innertopmargin=6pt,innerbottommargin=6pt]
\textbf{\small AT A GLANCE}\par\vspace{4pt}
\begin{description}[leftmargin=0pt,labelsep=4pt,itemsep=2pt,parsep=0pt,
    font=\small\bfseries]
  \item[Problem:] Extreme LLM quantization forces a choice between
    codebook quality (with lookup overhead) or binary speed (with
    poor accuracy).
  \item[Why it matters:] On-device deployment of large models is
    blocked by memory bandwidth; sub-1-bit compression at competitive
    quality would unlock a new class of edge applications.
  \item[Method:] Binary spherical coding maps weight blocks onto a
    pre-defined tight frame of $\pm1$ vectors; decoded via
    XNOR-popcount, requiring no lookup table.
  \item[Result:] LLaMA-3-8B at 1-bit storage within 1.6 perplexity
    points of 2-bit AQLM; 12$\times$ memory reduction versus FP16.
  \item[Evidence:] Benchmarked on WikiText-2 perplexity and MMLU
    accuracy for LLaMA-3 (8B, 13B, 70B) and Mistral-7B.
\end{description}
\end{mdframed}

\section{The Big Picture}

The deployment gap for large language models is primarily a memory
problem. A 70-billion-parameter model in FP16 requires $\sim$140\,GB
of GPU memory --- prohibitive for all but the largest datacenter
installations and completely out of reach for edge devices. Quantization
compresses weight values to fewer bits, reducing memory proportionally.

At 4-bit (GPTQ, AWQ) and 2-bit (AQLM, QuIP) precision, existing
methods achieve competitive accuracy, but even 2-bit precision requires
$\sim$17\,GB for a 70B model and a runtime codebook for reconstruction.
True 1-bit methods (BitNet, OneBit) eliminate the codebook but degrade
quality sharply because $\pm1$ scalar representations cannot capture
the heterogeneous magnitude structure of LLM weights.

BiSCo-LLM breaks out of this binary by separating the \emph{structure}
of the code (binary on a sphere) from its \emph{expressive power}
(controlled by the number of frame vectors $K$).

\section{Why Existing Approaches Fall Short}

\textbf{Codebook methods} (AQLM, QuIP, VQ-LLM): Represent weight blocks
as indices into a learnable codebook. Quality is excellent --- codebooks
can approximate diverse weight distributions --- but the codebook itself
must be stored in memory (typically 65,536 entries $\times$ 32 floats
$= 2$\,MB per layer) and queried at every matrix-vector multiply via
a memory-bandwidth-intensive lookup.

\textbf{Binary methods} (BitNet, OneBit): Represent weights as single
$\pm1$ scalars (possibly scaled per row or column). Zero lookup overhead
and excellent throughput, but accuracy degrades sharply because all
weight columns of different norms are mapped to the same code alphabet.
The fundamental mismatch is that $\pm1$ cannot represent graded
magnitudes.

\textbf{Vector quantization without lookup}: No prior work achieves
this combination. The key missing insight is that binary \emph{combinations}
of vectors, rather than individual binary scalars, can approximate
rich high-dimensional distributions while still decoding via purely
binary arithmetic.

\section{Core Mechanism}

A binary spherical code represents each weight column $w \in \mathbb{R}^d$
as a sparse binary combination of pre-defined frame vectors
$\{b_k\}_{k=1}^K \subset \{+1,-1\}^d$ that form a tight frame:
\[
\hat{w} = \frac{1}{\sqrt{K}} \sum_{k=1}^K z_k b_k, \quad z_k \in \{0,1\}
\]

The frame property $\sum_k b_k b_k^\top = (K/d) I_d$ guarantees
that the reconstruction operator is an isometry in expectation,
distributing reconstruction error uniformly across all weight directions
regardless of the weight norm distribution.

At runtime, decoding $\hat{w}$ requires only XNOR-popcount operations
between $z$ and the frame matrix --- no lookup table, no floating-point
multiply at decode time.

\section{Technical Formulation}

The binary tight frame satisfies:
\[
\min_{\{b_k\} \subset \{+1,-1\}^d}
\left\|\sum_{k=1}^K b_k b_k^\top - \frac{K}{d}\,I_d\right\|_F^2
\]

This Grassmannian binary packing problem is solved once offline via
a greedy sequential algorithm. Given the fixed frame, per-weight
compression minimizes reconstruction error:
\[
z^* = \operatornamewithlimits{argmin}_{z \in \{0,1\}^K}
\left\|w - \frac{1}{\sqrt{K}}\sum_{k=1}^K z_k b_k\right\|^2
\]

During training with straight-through estimators:
\[
\frac{\partial \mathcal{L}}{\partial z_k}
\approx \frac{\partial \mathcal{L}}{\partial \hat{w}} \cdot \frac{b_k}{\sqrt{K}}
\cdot \mathbf{1}\!\left[|z_k - 0.5| < 0.5\right]
\]

The effective bit-rate is $R = (\log_2 K)/d$ bits per weight,
tunable from $\approx 0.5$ to $\approx 1.5$ bits by varying $K \in [2^{d/3}, 2^{2d}]$.

\section{Learning or Inference Procedure}

\begin{enumerate}[itemsep=2pt,parsep=0pt]
  \item \textbf{Frame construction (once):} Build a binary tight frame
    $\{b_k\}_{k=1}^K$ offline via greedy orthogonal pursuit.
  \item \textbf{Initialization:} For each weight block $W$, solve the
    binary coding problem via fast iterative bit-flip optimization.
  \item \textbf{Fine-tuning:} Train the binary codes $z$ (not the frame)
    using STE gradients with a small calibration dataset ($\sim$512 sequences).
  \item \textbf{Inference:} Replace each weight block with its binary code.
    Matrix-vector products computed as $\hat{W}x \approx \frac{1}{\sqrt{K}} B_z x$
    where $B_z x$ is computed via XNOR-popcount on binary hardware.
\end{enumerate}

\section{What the Guarantee Says}

For a $d$-dimensional weight vector and a random binary tight frame
with $K$ vectors, the expected reconstruction error is:
\[
\mathbb{E}\!\left[\|w - \hat{w}\|^2\right]
= \|w\|^2 \cdot \left(1 - \frac{K}{d \cdot 2^K}\sum_{s=0}^K\binom{K}{s}
\left|\frac{s}{K}-0.5\right|\right)
\]

This bound is tight and shows that reconstruction error vanishes as
$K \to \infty$ (more frame vectors $=$ higher quality), independently
of the weight norm distribution.

\section{Experimental Findings}

\begin{center}
\small
\begin{tabular}{lcc}
\toprule
\textbf{Method (LLaMA-3-8B)} & \textbf{Bits} & \textbf{PPL} \\
\midrule
FP16            & 16  & 6.14 \\
GPTQ-4bit       & 4   & 6.58 \\
AQLM-2bit       & 2   & 7.21 \\
BitNet-1bit     & 1   & 9.87 \\
BiSCo-LLM 1bit  & 1   & 7.73 \\
BiSCo-LLM 0.5bit & 0.5 & 9.14 \\
\bottomrule
\end{tabular}
\end{center}

WikiText-2 perplexity (lower is better). BiSCo-LLM at 1-bit matches
AQLM's 2-bit quality within 0.5 PPL at half the memory footprint,
and outperforms BitNet by 2.1 PPL at the same storage budget.

\section{Ablations and Interpretation}

\textbf{Frame size $K$:} Quality improves monotonically with $K$ up to
$K \approx 2d$, after which reconstruction error saturates. The sweet
spot for 1-bit effective compression is $K = d$: each weight column uses
$d$ binary bits to select $d$ frame vectors, yielding exactly 1 bit/weight.

\textbf{Attention vs.\ MLP layers:} Attention weights are more compressible
(quality gap vs.\ FP16 smaller by 15\%) than MLP weights, consistent with
prior quantization literature. BiSCo-LLM supports layer-specific $K$ tuning.

\textbf{XNOR-popcount throughput:} On a Snapdragon X Elite mobile SoC,
BiSCo-LLM achieves $3.4\times$ higher token generation throughput versus
INT4 GPTQ, demonstrating that the lookup-free design translates directly
to real hardware efficiency.

\textbf{Model scale:} Perplexity gap vs.\ FP16 decreases at larger
model sizes (70B model: 1.8 PPL gap vs.\ 1.6 PPL gap for 8B at 1-bit),
suggesting BiSCo-LLM scales favorably.

\section*{Reference}

Yuantian Shao, Peisong Wang, Zhilei Liu, Chuangyi Li, Yuanteng Chen,
Pengcheng Xie, Yiwu Yao, Zhihui Wei, Jian Cheng.
\textbf{BiSCo-LLM: Lookup-Free Binary Spherical Coding for Extreme
Low-Bit Large Language Model Compression}.
arXiv:2607.08643, July 2026.
\url{https://arxiv.org/abs/2607.08643}

\end{document}
